\scanpage{37}
$\varprojlim_n({}_nT)_{\bar s_1}=({}_nG_m)^r_{\bar s_1}$, donc les opérations de $\pi_1(\widetilde S,\bar s_0)$ sur $M_\infty$ se font par l'intermédiaire du ``caractère \unsure{standart|standard}'' l'opération sur les racines de l'unité. Or si on a $\rho_{ab}\neq 0$, alors les opérations du gpe de Galois $\pi_{\deleted{\unsure{0}}1}(\widetilde S,\bar s_0)=\pi_{\unsure{1|0}}(s_0,\bar s_0)$ \uline{ne se font pas} par l'intermédiaire exclusif de ce caractère \ldots

On est ramené à un énoncé sur une V.A. sur un corps \deleted{\unsure{f|$\mathfrak{k}$}} de type fini [fini serait ici suffisant \ldots]

\begin{center}\rule{2.5cm}{0.35pt}\end{center}

Restent à prouver $\text{\textcircled{\scriptsize 3}}$ et $\text{\textcircled{\scriptsize\unsure{4|9}}}$.
\margin{\unsure{commencer} par cela ....}
% slanted

$\text{\textcircled{\scriptsize 3}}\Longleftrightarrow\text{\textcircled{\scriptsize 9}}$, et \deleted{\unsure{résulte}} \unsure{résulterait|résulte|résultat} du fait qu'\uline{un gpe résoluble se spécialise en un résoluble}

$\text{\textcircled{\scriptsize 4}}$ On regarde un tore maximal $T_0$ de $G_0$
\diagram{\begin{tikzcd}[ampersand replacement=\&,row sep=0pt,column sep=1em]
G\&G_1\\
\cup\&\cup\\
\overline{R_1}=R_0\&R_1
\end{tikzcd}}
\margin{$G_0/R_0$}
on prolonge les $ {}_nT$ en des $H(n)$ de façon compatible, alors on voit \unsure{(si|si|$G_1$} $r_0=\rho_{ss}(G_0)\inserted{\leq\dim T_0/T_0\cap R_0}$) que les $H(n)_1/R_1\cap H(n)_1$ sont de ``$n$-rang'' $\geq\dim T_0/T_0\cap R_0\geq r_0$. Ils sont de $n$-rang $\leq\rho'_{ss}$, d'où $\rho'_{ss}\geq\rho_{ss}$ \unsure{\textendash|\ldots}
% archive sheet number: 36
